\documentclass[10pt,a4paper]{article}
\usepackage[margin=1in]{geometry}
\usepackage{hyperref}
\hypersetup{
    colorlinks=true,
    linkcolor=blue,
    urlcolor=blue,
    citecolor=blue
}
\usepackage{enumitem}
\usepackage{titlesec}
\usepackage{parskip}
\usepackage{fontawesome5}  

\titleformat{\section}{\large\bfseries}{}{0em}{}[\titlerule]

\begin{document}

\begin{center}
    {\Huge \textbf{Fernando Gutiérrez Canales}}\\
    {\Large Data Scientist | PhD}\\    
    León, Mexico \quad | \quad \href{https://www.linkedin.com/in/fernando-guti%C3%A9rrez-canales-804876343/}{LinkedIn} \quad | \quad +52 1 477 55 18 832 \quad | \quad \href{mailto:carl.cfgc@gmail.com}{carl.cfgc@gmail.com} \\
     \href{https://github.com/Fernando-Canales}{GitHub: github.com/Fernando-Canales}
\end{center}

\vspace{1em}

\section*{Professional Summary}
Data Scientist with a PhD in Astrophysics and 5+ years of experience designing scalable data pipelines, evaluating Large Language Models (LLMs), and applying advanced statistical analysis to massive, complex datasets. Expert in Python, C, and machine learning frameworks. Proven ability to translate rigorous scientific research into optimized algorithmic solutions, bridging the gap between highly technical data architecture and impactful business outcomes.

\section*{Professional Experience}
\textbf{Mercor} \hfill \textbf{Göttingen (Remote)} \\
\textit{AI Prompt Engineer (LLM Training \& Evaluation)} \hfill \textit{Oct 2025 -- Feb 2026}
\begin{itemize}[leftmargin=1.5em]
	\item Developed complex logical and analytical datasets (adversarial prompts) to train and fine-tune Large Language Models (LLMs), effectively testing and improving their reasoning capacity by 30\%.
\end{itemize}

\textbf{Paris Observatory \& Max Planck Institute for Solar System Research} \hfill \textbf{Paris and Göttingen} \\
\textit{Data Scientist \& Researcher} \hfill \textit{Mar 2022 -- Jul 2025}
\begin{itemize}[leftmargin=1.5em]
    \item Architected and deployed automated Python data pipelines to process massive astronomical datasets, utilizing advanced statistical methods and time series analysis to estimate exoplanet detection efficiencies.
    \item Integrated C and Bash modules to optimize algorithmic performance, drastically reducing computational overhead and manual data handling time.
    \item Collaborated with scientists across 15+ international institutions to standardize complex data architectures and ensure stringent data integrity for the PLATO space mission.
\end{itemize}

\textbf{ESTEC, European Space Agency (ESA)} \hfill \textbf{Noordwijk, Netherlands}\\
\textit{Data Scientist (Internship)}  \hfill \textit{Jun 2023 -- Aug 2023}
\begin{itemize}[leftmargin=1.5em]
    \item Engineered automated scripts using Python to extract, validate, and statistically analyze complex parameters from photometric detectors, enhancing the reliability of mission-critical data pipelines.
\end{itemize}

\section*{Featured Data Science Projects}
\textbf{Predictive Modeling \& Statistical Analysis (TripleTen)} \hfill \textbf{Python, Scikit-learn, Pandas} \\
\textit{Data Scientist} \hfill \textit{March 2026 -- April 2026}
\begin{itemize}[leftmargin=1.5em]
    \item Developed a machine learning model using Scikit-learn to predict customer churn, achieving an accuracy of 85\%
    \item Conducted rigorous A/B testing and statistical hypothesis testing to validate the model's impact on business retention strategies
\end{itemize}

\textbf{Ophthalmology Clinic Data Architecture} \hfill \textbf{León, Mexico} \\
\textit{Lead Data Consultant (Independent Project)} \hfill \textit{2021 -- 2022}
\begin{itemize}[leftmargin=1.5em]
    \item Executed end-to-end data wrangling on years of unstructured physical and digital records, building a centralized, robust relational database.
    \item Designed automated scripts to extract actionable insights, reducing information retrieval time and optimizing clinic resource allocation.
\end{itemize}

\section*{Technical Skills}
\textbf{Programming \& Core Stack:} Python (Pandas, NumPy, Scikit-learn, SciPy), C, Bash, SQL, R. \\
\textbf{Data Science \& AI:} Machine Learning, LLM Prompt Engineering, Statistical Modeling, Time Series Analysis, Predictive Analytics, Algorithm Optimization. \\
\textbf{Data Engineering \& Tools:} Automated Pipelines, Data Architecture, Git, Docker, Linux, Tableau.

\section*{Education}
\textbf{Data Analytics Bootcamp Certification} \hfill \textit{TripleTen} \hfill \textit{2026}\\
\textbf{PhD in Astrophysics} \hfill \textit{Paris Observatory \& MPS Göttingen} \hfill \textit{2025}\\
\textbf{Master's Degree in Astrophysics} \hfill \textit{University of Guanajuato, Mexico} \hfill \textit{2021}\\
\textbf{Bachelor's Degree in Physics} \hfill \textit{University of Guanajuato, Mexico} \hfill \textit{2019}

\section*{Achievements \& Publications}
\begin{itemize}[leftmargin=1.5em]
    \item Co-author of \href{https://ui.adsabs.harvard.edu/search/fq...}{peer-reviewed publications} leveraging advanced statistical methods for photometric data analysis.
    \item Selected to present predictive research findings at major international conferences (EAS Padua 2024, PLATO Week Prague 2023).
    \item Recipient of full academic scholarships and graduated with Honorable Mention across all three university degrees (2019, 2021, 2025).
\end{itemize}

\end{document}
